\documentclass[11pt]{article}

\usepackage[top=1in, bottom=1.2in, left=0.85in, right=0.7in]{geometry}
\usepackage{authblk}
\usepackage{hyperref}
\usepackage[utf8]{inputenc}
\usepackage{amsmath}
\usepackage{booktabs}
\usepackage{amsfonts}
\usepackage{amssymb}
\usepackage{siunitx}
\usepackage{graphicx}
\usepackage{subcaption}
\usepackage{float}
\usepackage{indentfirst}
\usepackage{longtable}
\usepackage{makecell} % Add this in your preamble
\usepackage{booktabs}
\usepackage{threeparttable}
\usepackage{setspace}
\usepackage{natbib}
\newcommand{\sym}[1]{\ifmmode^{#1}\else\(^{#1}\)\fi}
\usepackage[margin=1in]{geometry}

\usepackage[backend=biber, style=numeric, sorting=unsrt]{biblatex}
\addbibresource{References.bib} % Ensure your .bib file name matches

\onehalfspacing

\title{Computational Economics: Problem Set 8}
\author{Emma (Sunwoo) Cho}
\date{\today}  % or leave empty: \date{}

\begin{document}

\maketitle

\section{Economic Model}

Following \cite{CooperEjarque2003}, the firm chooses its capital stock in the presence of market power and convex adjustment costs. The value function $V(K,A)$ solves
\begin{equation}
  V(K,A)
  =
  \max_{K'}
  \left\{
    \pi(K,A)
    - p\bigl(K' - (1-\delta)K\bigr)
    - C(K',K)
    + \beta\,\mathbb{E}_{A'|A}\bigl[ V(K',A') \bigr]
  \right\}.
  \label{eq:bellman_basic}
\end{equation}

for all $(K,A)$, where $K$ is the current capital stock and $A$ is a measure of profitability. 

%In a microfoundation based on a downward-sloping demand curve and a Cobb–Douglas technology, $A$ bundles shocks to productivity, factor prices and demand into a reduced-form profitability term driving $\pi(K,A)$.} 

Profits are given as follow
\begin{equation}
  \pi(K,A) = A K^{\alpha},
\end{equation}
where $\alpha$ measures the degree of market power. 
Capital evolves according to $K' = (1-\delta)K + I$, and adjustment costs are quadratic in investment relative to the capital stock:
\begin{equation}
  C(K',K)
  =
  \frac{\gamma}{2}
  \left( \frac{K' - (1-\delta)K}{K} \right)^{2} K.
\end{equation}

The profitability shock $A$ follows a stochastic process governed by parameters $(\rho,\sigma)$. \cite{CooperEjarque2003} treat $(\rho,\sigma)$ as parameters of the process for the firm-specific shocks to $A$, and approximates this process using a discrete state space constructed by the method of \cite{Tauchen1986}. 
In the PS8 instructions, this is implemented as an AR(1) in logs, $\ln A_{t+1} = \rho \ln A_t + \varepsilon_{t+1}$ with $\varepsilon_{t+1} \sim \mathcal{N}(0,\sigma^2)$, before discretization.(Followed by PS8 instruction, which specifies that firm productivity follows an AR(1) process in logs and that you should discretise it using the Tauchen (1986) method.)

\subsection{Costly External Finance}

Section 3.2 of \cite{CooperEjarque2003} introduces costly external finance. The cost of accessing capital markets is
\begin{equation}
  \phi(E) = \phi_0 + \phi_1 E,
\end{equation}
where $E$ measures the extent of external finance and is defined by
\begin{equation}
  E \equiv I - \pi(A,K).
\end{equation}

With this cost of capital-market participation, the firm’s dynamic problem becomes
\begin{equation}
  V(K,A) = \max\bigl\{ V^{e}(K,A),\, V^{i}(K,A) \bigr\},
  \label{eq:bellman_split}
\end{equation}
where $V^{e}$ and $V^{i}$ are the value functions under external and internal finance, respectively. \cite{CooperEjarque2003} write the external-finance problem as
\begin{align}
  V^{e}(K,A)
  =
  \max_{K' > \pi(K,A) + (1-\delta)K}
  \Bigl\{
    &\pi(K,A)
    - p\bigl(K' - (1-\delta)K\bigr)
    - C(K',K)
    - \phi_0
    - \phi_1 \, p\bigl(K' - K\bigr(1-\delta)). \nonumber \\
    &\left. - \pi(K,A)\bigr)
    + \beta\,\mathbb{E}_{A'|A}\bigl[ V(K',A') \bigr]
  \Bigr\},
  \label{eq:Ve}
\end{align}
and the internal-finance problem as
\begin{align}
  V^{i}(K,A)
  =
  \max_{K' \leq \pi(K,A) + (1-\delta)K}
  \Bigl\{
    \pi(K,A)
    - p\bigl(K' - K\bigr(1-\delta))
    - C(K',K)
    + \beta\,\mathbb{E}_{A'|A}\bigl[ V(K',A') \bigr]
  \Bigr\}
  \label{eq:Vi}
\end{align}
for all $(K,A)$.

For the specification in Table~3, the authors define the fixed cost relative to mean capital, $\tilde{\phi}_0$, and estimate the parameter vector
\[
  \Theta = (\alpha,\gamma,\rho,\sigma,\tilde{\phi}_0),
\]
and set to $\phi_1 = 0$.

\section{Estimation Strategy}

The problem set report's goal (objective) is to replicate the costly external finance part of Table~3 using a Simulated Method of Moments (SMM) estimator. Let $\hat{\psi}^d$ denote the vector of data moments reported in Table~3(b) and $\hat{\psi}^s(\theta)$ the corresponding moments computed from simulated panels generated by the model. The SMM estimator solves
\begin{equation}
  \hat{\theta}_{S,T}(W)
  =
  \arg\min_{\theta}
  \left[
    \hat{\psi}^d
    -
    \hat{\psi}^s(\theta)
  \right]^{\!\top}
  W
  \left[
    \hat{\psi}^d
    -
    \hat{\psi}^s(\theta)
  \right],
  \label{eq:smm_criterion}
\end{equation}
where $W$ is a positive definite weighting matrix.

PS8 instruction specifies that first estimate the model with $W$ equal to the identity matrix, and then construct an optimal weighting matrix from the variance–covariance matrix of the simulated moments using multiple simulations at the first-stage estimate. 

\section{Results}

The estimation results are presented in Table \ref{tab:results}. The estimated fixed cost of external finance, $\tilde{\phi}_0$, is statistically indistinguishable from zero from this simulation experiment.

\begin{table}[htbp]
\centering
\caption{SMM Estimation Results: Replication of Cooper \& Ejarque (2003) Table 3}
\label{tab:results}
\vspace{0.2cm}
\textbf{Panel A: Structural Parameter Estimates} \\
\vspace{0.1cm}
\begin{tabular}{lccc|c}
\hline\hline
& \multicolumn{2}{c}{\textbf{Our Replication}} & & \textbf{Original Paper} \\
Parameter & Estimate & Std. Err. & & Estimate (Table 3) \\
\hline
Curvature ($\alpha$)          & 0.7411 & (0.0006) & & 0.696 \\
Adjustment Cost ($\gamma$)    & 0.3518 & (0.0019) & & 0.133 \\
Persistence ($\rho$)          & 0.6783 & (0.0016) & & 0.893 \\
Shock Volatility ($\sigma$)   & 0.2478 & (0.0007) & & 0.098 \\
Fixed Fin. Cost ($\tilde{\phi}_0$) & 0.0005 & (0.0005) & & 0.000 \\
\hline\hline
\end{tabular}

\vspace{0.5cm}
\textbf{Panel B: Moments (Data vs. Model)} \\
\vspace{0.1cm}
\begin{tabular}{lcc}
\hline\hline
Moment & Data Target & Model Fit \\
\hline
Regression coeff on $Q$ ($a_1$)       & 0.030 & 0.125 \\
Regression coeff on Cash Flow ($a_2$) & 0.240 & 1.162 \\
Serial Correlation of $I/K$           & 0.400 & 0.375 \\
Std. Dev. of Profit Rate ($\pi/K$)    & 0.250 & 0.091 \\
Mean Average $Q$ ($\bar{q}$)          & 3.000 & 2.819 \\
External Finance Fraction             & 0.250 & 0.079 \\
\hline
$J$-Statistic & \multicolumn{2}{c}{31140.44} \\
\hline\hline
\end{tabular}
\end{table}




\begin{enumerate}
    The estimation results yields $\hat{\tilde{\phi}}_0 = 0.0005$ with a standard error of $0.0005$. This is statistically indistinguishable from zero and very close to the paper's estimate of $0$. This confirms the paper's main conclusion, which was that once market power is accounted for, costly external finance is not required to match investment moments.
    
    \item $\hat{\alpha} \approx 0.74$, which is close to the paper's estimate of $0.696$. This confirms the presence of diminishing returns to scale (market power), which is the primary driver of the gap between average $Q$ and marginal $Q$.
    
    \item While replicating, it ended up with a slightly different structural combination than the authors. The simulation code result shows that higher adjustment costs ($\hat{\gamma} \approx 0.35$ vs $0.13$) and lower persistence ($\hat{\rho} \approx 0.68$ vs $0.89$). It relied more heavily on adjustment costs ($\gamma$) to match the serial correlation target ($0.375$ vs. target $0.40$), whereas the authors relied more on shock persistence (this part should be improved to replicate the results similarly further).
\end{enumerate}

\subsection{Numerical Details}

The dynamic programming problem is solved using Value Function Iteration (VFI) on a discretised state space. 

\begin{itemize}
    \item The continuous productivity process $\ln A_t$ is approximated by a Markov chain with $N_A = 7$ states using the method of \cite{Tauchen1986}. The capital stock $K_t$ is discretized into a grid of $N_K = 60$ points linearly spaced between $K_{\min}=1.0$ and $K_{\max}=350.0$.
    \item The VFI algorithm terminates when the maximum absolute difference between successive value functions falls below a tolerance of $1 \times 10^{-4}$.
    \item For this experiment, I simulated a panel of $N=500$ firms for $T=30$ periods. Continuous uniform random draws are mapped to the discrete productivity states using inverse transform sampling on the Tauchen transition matrix.
    \item Nelder-Mead simplex algorithm has been employed for the minimisation because the discretisation of the state space via grids leads to a non-smooth SMM objective function. 
\end{itemize}
\printbibliography
\end{document}


